\documentclass[12pt]{article}
\usepackage[T1]{fontenc}
\usepackage[utf8]{inputenc}
\usepackage{lmodern}
\usepackage{textcomp}
\usepackage{enumitem}
\usepackage{geometry}
\usepackage{graphicx}
\usepackage{amsmath, amssymb}
\geometry{margin=1in}
\begin{document}
\begin{center}
Sociology - Undergraduate Level Assessment 11 \
Total Marks: 40 \
Answer all questions.
\end{center}

\section*{Multiple Choice (10 marks)}
\begin{enumerate}[label=\arabic*.]
\item Sullivan's (2000) study suggested that the proportion of housework men did was greatest when:
\begin{enumerate}[label=(\alph*)]
    \item they had rediscovered themselves as 'new men'
    \item their wives were at home and nagged them all the time
    \item exciting gadgets like the hoover and electric iron were invented
    \item they were unemployed or both partners worked full time
\end{enumerate}
\item In the context of the labour movement in the nineteenth century, 'incorporation' meant:
\begin{enumerate}[label=(\alph*)]
    \item including union representatives in processes of policy decision making
    \item creating links between the state and corporate organizations
    \item recruiting women into full time paid employment
    \item including working class organizations in political bargaining and representation
\end{enumerate}
\item Durkheim defined social facts as:
\begin{enumerate}[label=(\alph*)]
    \item ways of acting, thinking, and feeling that are collective and social in origin
    \item the way scientists construct knowledge in a social context
    \item data collected about social phenomena that are proven to be correct
    \item ideas and theories that have no basis in the external, physical world
\end{enumerate}
\item Marx said that the development of the labour movement through factory-based production would turn the working class into:
\begin{enumerate}[label=(\alph*)]
    \item a class in itself
    \item a class by itself
    \item a class for itself
    \item a ruling class
\end{enumerate}
\item Society cannot be studied in the same way as the natural world because:
\begin{enumerate}[label=(\alph*)]
    \item human behaviour is meaningful, and varies between individuals and cultures
    \item it is difficult for sociologists to gain access to a research laboratory
    \item sociologists are not rational or critical enough in their approach
    \item we cannot collect empirical data about social life
\end{enumerate}
\end{enumerate}

\section*{True / False (10 marks)}
\textit{Answer True or False}
\begin{enumerate}[label=\arabic*.]
\setcounter{enumi}{5}
\item True or False: The experimental method relies on personal beliefs and values to determine what research to conduct.
\item True or False: Totalitarian societies typically restrict freedom of movement for their citizens to maintain control.
\item True or False: A fixed geographical location is not considered a necessary criterion of community by Fulcher \& Scott.
\item True or False: The human relations approach emphasizes teamwork, effective communication, and employee satisfaction in the workplace.
\item True or False: Sociology relies mainly on subjective interpretations rather than systematic data collection and empirical evidence.
\end{enumerate}

\section*{Long Form Response (20 marks)}
\textit{Answer each question with a clear and complete explanation.}
\begin{enumerate}[label=\arabic*.]
\setcounter{enumi}{10}
\item Discuss the characteristics and societal implications of the deviant subculture of homosexuals known as "mollies" in seventeenth and eighteenth century London.
\item Discuss the features of globalization and critically analyze how the concept challenges the traditional power of nation-states in contemporary society.
\end{enumerate}

\vfill
\noindent\textit{End of Paper}
\end{document}